\section{动机}
\label{sec:motivation}

\subsection{挑战}
当小型本地模型主导调查时，会出现三个挑战。我们在各项研究的基线配置（\S\ref{sec:setup}）中逐一测量了它们；在这些配置里，模型自己选择探针。

\stepnum{1} \textbf{探测却不收敛。}自己选择探针的 Qwen2.5-7B 只验证了 0.083 到 0.125 的案例。它每个回合运行约七个不同的探针，花掉 10 个成本单位中的 9.4 到 9.5 个，在 72 个回合里只有 6 或 7 个给出了结论；四分之三的回合以工具预算耗尽结束。把它的选择换成均匀随机的选择，完成率提高 $+0.556$（\S\ref{sec:rq-rank}），说明它自己的选择比随机还差。

\stepnum{2} \textbf{从不决定停止。}Llama-3.1-8B 在全部 3{,}389 次决策中每次都要求再运行一个探针，包括账本后验已超过 0.999、且提示中已说明没有合法探针可用的时候。它的输出始终是合法的 JSON；它只是从不结束。更好的探测帮不了一个从不给出结论的模型。

\stepnum{3} \textbf{没有支撑的结论。}更大的模型会下结论，但并不总是基于证据。在没有账本的情况下，Qwen2.5-32B 和 72B 把 25\% 和 37\% 的需处置事件判为良性，引用中分别有 22\% 和 31\% 是无效证据（\S\ref{sec:security}）。在分诊中，把真实攻击判为无害的结论代价最高，因为事件会在未作响应的情况下被关闭。

\subsection{我们的方案}
下面说明 \hesp{} 如何应对这些挑战。

\stepnum{1} \textbf{在模型之外做信息增益选择。}\hesp{} 按每单位成本对候选起因不确定性的期望降低量，为每个合法探针排序，并执行排名第一的探针，不论模型提议的是什么。排序所依据的似然值从无 LLM 的开发运行中计数得到，而不是让模型给出——模型给不出来：模型给出的表使完成率只有 0.069 到 0.375（\S\ref{sec:rq-help}）。

\stepnum{2} \textbf{控制器侧停止。}在每次调用规划器之前，\hesp{} 应用一条明确的接受规则。一旦领先起因的后验至少为 0.8、且有一条当前观测支持它，控制器就自行下结论并引用该证据。模型仍可以更早结束，因此强模型保留了等待确认证据的选择。

\stepnum{3} \textbf{受证据约束的结论。}\hesp{} 只接受引用了当前观测、且该观测提高了所声称起因分数的结论。一个规划器无法访问的独立验证器，会用只有该起因才会产生的特征来检查每个结论；预写日志使审计能够重放每一步。
